\documentclass[10pt]{article}

\RequirePackage[top=0.8in,left=0.8in,right=0.8in,bottom=0.3in,includefoot]{geometry}
 

% Include a series of generic definitions and usepackage commands
% mediabuilder.tex is an external preamble that can be included in 
% pandoc templates to improve the way tex templates are rendered
% This file contains a bunch of custom includes that I use across many tex
% templates

% Stuff added for pandoc

% ALGORITHMS -------------------------------------------------------------------

% These two packages are used to write algorithms
\usepackage{algorithm}
\usepackage{algpseudocode}


% HYPERLINKS -------------------------------------------------------------------

% Setting up hyperref so that links can be populated
\usepackage{ifxetex,ifluatex}
\ifxetex
  \usepackage[setpagesize=false, % page size defined by xetex
              unicode=false, % unicode breaks when used with xetex
              xetex]{hyperref}
\else
  \usepackage[unicode=true]{hyperref}
\fi
\hypersetup{breaklinks=true, pdfborder={0 0 0}}


% FIGURES ----------------------------------------------------------------------
% used for wrapfig
\usepackage{wrapfig}

\usepackage{dblfloatfix}

% To handle multicolumn figures
% This lets you do \Begin{widefig}
\let\Begin\begin
\let\End\end
\newenvironment{widefig}{
	\renewenvironment{figure}{\begin{figure*}\centering}{\end{figure*}}
	\renewenvironment{table}{\begin{table*}\centering}{\end{table*}}
}

\newenvironment{widefigbot}{
  \renewenvironment{figure}{\begin{figure*}[!b]\centering}{\end{figure*}}
  \renewenvironment{table}{\begin{table*}\centering}{\end{table*}}
}

% Scale images if necessary, so that they will not overflow the page
% margins by default, and it is still possible to overwrite the defaults
% using explicit options in \includegraphics[width, height, ...]{}
\usepackage{graphicx}
\makeatletter
\def\maxwidth{\ifdim\Gin@nat@width>\linewidth\linewidth\else\Gin@nat@width\fi}
\def\maxheight{\ifdim\Gin@nat@height>\textheight\textheight\else\Gin@nat@height\fi}
\makeatother
\setkeys{Gin}{width=\maxwidth,height=\maxheight,keepaspectratio}

% Old way to do this:
% $if(graphics)$
% \usepackage{graphicx}
% % We will generate all images so they have a width \maxwidth. This means
% % that they will get their normal width if they fit onto the page, but
% % are scaled down if they would overflow the margins.
% \makeatletter
% \def\maxwidth{\ifdim\Gin@nat@width>\linewidth\linewidth
% \else\Gin@nat@width\fi}
% \makeatother
% \let\Oldincludegraphics\includegraphics
% \renewcommand{\includegraphics}[1]{\Oldincludegraphics[width=\maxwidth]{#1}}
% $endif$


% When I was using the multicol environment, I had to redefine the figure
% environment so there's no float (what pandoc does by default)
% to make it compatible with the multicol package. This is how I did that.
% This led to several problems, such as figures showing up out-of-order. In the
% end, because of this, I abandoned the multicol package and switched back to
% the basic, builtin twocol environment. Leaving this here in case I ever need
% multicol again:
% \usepackage{float}
% \let\origfigure\figure
% \let\endorigfigure\endfigure
% \renewenvironment{figure}[1][2] {
%     \expandafter\origfigure\expandafter[H]
% } {
%     \endorigfigure
% }

% % Also for the table
% \let\origtable\table
% \let\endorigtable\endtable
% \renewenvironment{table}[1][2] {
%     \expandafter\origtable\expandafter[H]
% } {
%     \endorigtable
% }





% LISTS ----------------------------------------------------------------------

\let\tightlist\relax

% \providecommand{\tightlist}{
%   \setlength{\itemsep}{0pt}
%   \setlength{\parskip}{0pt}}

\providecommand{\tightlist}{
  \setlength{\topsep}{1em}
  \setlength{\partopsep}{1em}
  \setlength{\parskip}{1em}
  \setlength{\parsep}{1em}
  \setlength{\parindent}{0pt}
  \setlength{\itemsep}{0pt}
  \setlength{\itemindent}{0pt}
  \setlength{\leftmargin}{0pt}
}

% make bullets smaller
\def\labelitemi{\raisebox{0.5ex}{\tiny\textbullet}}
% this itemindent code adjusts the indentation of bullet lists.

% Redefine labelwidth for lists; otherwise, the enumerate package will cause
% markers to extend beyond the left margin.
\makeatletter\AtBeginDocument{%
  \renewcommand{\@listi}
    {\setlength{\labelwidth}{4em}}
}\makeatother
\usepackage{enumerate}



% TABLES -----------------------------------------------------------------------
% longtable is required for markdown-format tables to be rendered in latex
\usepackage{longtable}

% LAYOUT -----------------------------------------------------------------------
\def\topfraction{0.9} % 90 percent of the page may be used by floats on top
\def\bottomfraction{0.9} % 90 percent of the page may be used by floats at the bottom

\clubpenalty=9996
\widowpenalty=9999
\brokenpenalty=4991
\predisplaypenalty=10000
\postdisplaypenalty=1549
\displaywidowpenalty=1602

\setlength{\parindent}{0pt}
\setlength{\parskip}{6pt plus 2pt minus 1pt}
\setlength{\emergencystretch}{3em}  % prevent overfull lines


% FONTS ---

\renewcommand{\rmdefault}{phv} % Arial
\renewcommand{\sfdefault}{phv} % Arial
% \renewcommand{\rmdefault}{bch} % Charter
% \renewcommand{\sfdefault}{bch} % Charter




% \usepackage{microtype}

% Need this for the \Letter glyph, envelope, for corresponding author.
\usepackage[misc,geometry]{ifsym} 


\usepackage{authblk}

\renewcommand{\rmdefault}{bch} % Charter
\renewcommand{\sfdefault}{bch} % Charter


% used for twocol flush balance end on last page
\usepackage{flushend}
\usepackage{pdfpages}


\usepackage[skip=2pt, font=footnotesize, format=plain, labelfont=it, textfont=it, tableposition=below]{caption}
%  This code to use {caption} will wipe out the ability of pandoc-crossref
%  to prepend the "Figure" reference to the caption.
\DeclareCaptionLabelFormat{nolabel}{}
\captionsetup{labelformat=nolabel}
\usepackage{amssymb,amsmath}
\usepackage{ifxetex,ifluatex}

% colors
\definecolor{bluehaze}{RGB}{54,95,145}
\definecolor{maroonhaze}{RGB}{148,54,52}
\definecolor{orangehaze}{RGB}{227,108,10}
\definecolor{black}{RGB}{0,0,0}
\definecolor{bluetint}{RGB}{0,50,100}
\definecolor{bluetint2}{RGB}{0,30,60}
\definecolor{graytint}{RGB}{175,175,175}
\definecolor{graybit}{RGB}{245,245,245}
\definecolor{lightblack}{RGB}{75,75,75}
\definecolor{brickred}{rgb}{0.8, 0.25, 0.33}
%  This is used to allow Shaded environment


% Start column longtable hack
% Allows tables to be used in two-column documents with pandoc
% Credit: https://tex.stackexchange.com/a/224096/161489
\makeatletter
\let\oldlt\longtable
\let\endoldlt\endlongtable

\def\longtable{\@ifnextchar[\longtable@i \longtable@ii}
\def\longtable@i[#1]{\begin{figure}[H]
\onecolumn
\begin{minipage}{0.5\textwidth}
\oldlt[#1]
}
\def\longtable@ii{\begin{figure}[H]
\onecolumn
\begin{minipage}{0.5\textwidth}
\oldlt
}
\def\endlongtable{\endoldlt
\end{minipage}
\twocolumn
\end{figure}}
\makeatother
% End column longtable hack


\ifxetex
  \usepackage{fontspec,xltxtra,xunicode}
  \defaultfontfeatures{Mapping=tex-text,Scale=MatchLowercase}
\else
  \ifluatex
    \usepackage{fontspec}
    \defaultfontfeatures{Mapping=tex-text,Scale=MatchLowercase}
  \else
    \usepackage[utf8]{inputenc}
  \fi
\fi

\usepackage{ctable}
\usepackage{float} % provides the H option for float placement


\hypersetup{breaklinks=true, pdfborder={0 0 0}, colorlinks=true, urlcolor=bluehaze, linkcolor=bluehaze, citecolor=bluehaze, anchorcolor=bluehaze}
\setcounter{secnumdepth}{0}

\title{Example article title}
% % \author{true}
% 



\linespread{1.0}


% Fancy page numbering ---------------------------------------------------------

%\pagestyle{fancy}
%\fancyhead[RO]{Sheffield \thepage}
\usepackage{fancyhdr}
\fancyhead{} % empty
\fancyfoot{} % empty out the default, which is page number in center
\def\middot{\textperiodcentered~}
\lfoot{\scriptsize \thepage \middot Short title   }
\pagestyle{fancy}


% Fancy abstract box styling ---------------------------------------------------

\usepackage{mdframed}   % for framing
\newmdenv[linecolor=white, fontcolor=bluetint, backgroundcolor=graybit, leftmargin=8, rightmargin=8, innertopmargin=8, innerbottommargin=8, innerleftmargin=12, innerrightmargin=12, roundcorner=10pt, font=\normalsize]{infobox}



% Title spacing, size, and color ------------------------------------------------------
\usepackage{titlesec}
\usepackage{xcolor}% http://ctan.org/pkg/xcolor

\titleformat{\section}
  {\normalfont\fontfamily{phv}\fontsize{11pt}{11pt}
    \selectfont\color{bluetint}\bfseries}{\thesection}{1em}{}

\titleformat{\subsection}
  {\normalfont\fontfamily{phv}\fontsize{9pt}{9pt}
    \selectfont\color{black}\bfseries}{\thesubsection}{1em}{}

\titleformat{\subsubsection}
  {\normalfont\fontfamily{phv}\fontsize{1em}{1em}
    \selectfont\color{black}}{\thesubsubsection}{1em}{}



% use amsthm for addpunct, to add periods to the end of non-punctuated
% paragraph titles (header level 4; ####)
\usepackage{amsthm}
\makeatletter
\newcommand{\addperiod}[1]{#1\@addpunct{.}}
\makeatother
% use runin to specify no newline
\titleformat{\paragraph}[runin]
  {\normalfont\color{lightblack}\itshape}{\theparagraph}{}{\addperiod}

% \titlespacing{command}{left spacing}{before spacing}{after spacing}[right]
\titlespacing\section{0pt}{12pt plus 2pt minus 2pt}{-2pt plus 2pt minus 2pt}
\titlespacing\subsection{0pt}{11pt plus 2pt minus 2pt}{-2pt plus 2pt minus 2pt}
\titlespacing\subsubsection{0pt}{4pt plus 2pt minus 2pt}{-2pt plus 2pt minus 2pt}

\titlespacing{\paragraph}{0pt}{1pt plus 4pt minus 2pt}{6pt plus 2pt minus 2pt}


% The maximum part of the top or bottom of a text page that can be
% occupied by floats. This is set so that at least four lines of text 
% fit on the page.
% \renewcommand\topfraction{.9}
% \renewcommand\bottomfraction{.9}
% % The minimum amount of a text page that must be occupied by text.
% % This should accomodate four lines of text.
% \renewcommand\textfraction{.06}
% % The minimum amount of a float page that must be occupied by floats.
% \renewcommand\floatpagefraction{.94}

% % The same parameters repeated for double column output
% \renewcommand\dbltopfraction{.9}
% \renewcommand\dblfloatpagefraction{.9}

% Attempt to smash lists together

\let\tightlist\relax
\providecommand{\tightlist}{%
  \setlength{\itemsep}{0pt}
  \setlength{\parskip}{3pt}
  \setlength{\topsep}{-6pt}
  }
\usepackage{enumerate}
\usepackage{enumitem}
\setlist[itemize]{topsep=0pt,parsep=2pt,itemindent=0pt,leftmargin=18pt}
% \def\labelitemi{\raisebox{0.5ex}{\tiny\textbullet}}


% Convert author/affiliations from yaml into latex -----------------------------

      \author[1,2,3,4,\Letter]{Nathan C. Sheffield}
  
% institutions must come after \author definitions so that corresponding tag comes last.
\affil[1]{Center for Public Health Genomics, University of Virginia}
\affil[2]{Department of Public Health Sciences, University of Virginia}
\affil[3]{Department of Biomedical Engineering, University of Virginia}
\affil[4]{Department of Biochemistry and Molecular Genetics, University
of Virginia}

      \affil[ ]{\Letter \hspace{0.08cm}  Correspondence: \href{mailto:nsheffield@virginia.edu}{nsheffield@virginia.edu}}


% bring affils a little closer to author.
\setlength{\affilsep}{0.5em}

\usepackage{titling}
\renewcommand{\headrulewidth}{0pt}
\renewcommand{\Affilfont}{\scriptsize}
\makeatletter

% Define the nice-looking title for the databio theme
\renewcommand{\maketitle}{
      \begin{flushright}{\huge{\color{graytint}\textbf{\MakeUppercase{Research
Article}}}}\end{flushright}\vspace{-0.4em}
    {\color{graytint}\hrule}\vspace{0.3em}
    \begin{raggedright} 
  \huge\textbf{\@title}
    \par\normalsize \vspace{0.5em}
    {        \ignorespaces\@author
    
        \par}
  
     \vspace{0.2em}
    \end{raggedright}
} % end \maketitle
\makeatother

%\subsectionfont{\color{maroonhaze}}
%\subsubsectionfont{\color{orangehaze}}
% end colors

% this code uses titles.sty to make a pagebreak with every header 1
% so each section will begin on a new page.
%\newcommand{\sectionbreak}{\clearpage}



\newlength{\cslhangindent}       % set up new length
\setlength{\cslhangindent}{0em}
\newenvironment{cslreferences}%
  {\everypar{\setlength{\hangindent}{\cslhangindent}}}%
  {\par} % by default, this env does not change anything






\begin{document}

\twocolumn[{%
  \begin{@twocolumnfalse}
        \maketitle
    
        \begin{infobox}
    Vivamus eu rhoncus neque. Quisque egestas venenatis odio a mattis.
    Ut ligula turpis, facilisis a cursus eget, semper quis dolor.
    Integer varius est ipsum, porttitor ornare eros placerat eget. Nulla
    aliquet nisi arcu, sed vestibulum urna faucibus pretium. Maecenas
    laoreet diam non urna tincidunt iaculis a ut ex. Aenean sem enim,
    laoreet id accumsan sed, faucibus vitae diam. Aenean facilisis
    tincidunt risus. Mauris sit amet hendrerit est, sit amet maximus
    augue.
    \end{infobox}
    
          \end{@twocolumnfalse}
}]


\hypertarget{citations}{%
\section{Citations}\label{citations}}

Cite papers using brackets and \texttt{bibtex} keys. Example citation:
\texttt{{[}@Sheffield2016{]}} will be rendered like
this\textsuperscript{1}. Use semicolons to separate multiple
citations\textsuperscript{1,2}.

\hypertarget{figures}{%
\section{Figures}\label{figures}}

\begin{figure}
\centering
\includegraphics{fig/example_figure.pdf}
\caption{\label{abstract}Fig. \ref{abstract}: Example full-width figure}
\end{figure}

Refer to a figure using figure labels, so they are numbered
automatically, like this: \texttt{\textbackslash{}ref\{abstract\}} (See
Fig. \ref{abstract}). Wrap a figure using the \texttt{pandoc-wrapfig}
extension by adding `\{0\}' to the end of the caption (Fig.
\ref{wrapped}).

\setlength{\intextsep}{2pt}\setlength{\columnsep}{8pt}\begin{wrapfigure}{R}{44mm}\centering\includegraphics{fig/example_figure.pdf}\caption{\label{wrapped}Fig.
\ref{wrapped}: Example wrapped figure}\vspace{-5pt}\end{wrapfigure}

Duis in tempor mauris, a lobortis nisl. Integer arcu lorem, vehicula sed
ante commodo, maximus eleifend nisi. Aenean efficitur molestie lorem, ac
pharetra felis euismod nec. Duis vitae ligula facilisis, dignissim justo
eget, elementum est. Nulla quis mi a justo porta pellentesque eget sit
amet purus. Ut ac vestibulum ante, in efficitur massa. Cras feugiat in
urna facilisis ultrices. Nullam vestibulum, lacus eget pretium pharetra,
augue ligula consectetur diam, eget condimentum ipsum magna sed augue.

\setlength{\intextsep}{2pt}\setlength{\columnsep}{8pt}\begin{figure*}[t]\centering\includegraphics{fig/example_figure.pdf}\caption{\label{fig3}Fig.
\ref{fig3}: Example double-column figure}\vspace{-5pt}\end{figure*}

Lorem ipsum dolor sit amet, consectetur adipiscing elit. Sed vulputate
tellus sed vulputate maximus. Interdum et malesuada fames ac ante ipsum
primis in faucibus. Praesent interdum sapien dui, in placerat velit
egestas auctor. Class aptent taciti sociosqu ad litora torquent per
conubia nostra, per inceptos himenaeos. Etiam ac nunc auctor, iaculis
nulla quis, tempor arcu. Vestibulum volutpat, massa sed blandit aliquet,
est elit rhoncus est, eu finibus leo turpis sit amet nisl. Fusce nec
diam ut sapien rutrum consequat. Aliquam massa sapien, molestie sed
ligula vitae, lacinia rutrum mi. Fusce ullamcorper accumsan quam vel
aliquam. Nullam efficitur, arcu eu consectetur tempus, mi dui facilisis
dolor, id lacinia augue turpis ut massa. Proin tincidunt congue
lobortis. Vestibulum semper condimentum enim accumsan malesuada. Aenean
interdum tellus id ex scelerisque mattis. Mauris tempor urna felis, eu
efficitur eros tempus condimentum. Aliquam erat volutpat.

Vivamus eu rhoncus neque. Quisque egestas venenatis odio a mattis. Ut
ligula turpis, facilisis a cursus eget, semper quis dolor. Integer
varius est ipsum, porttitor ornare eros placerat eget. Nulla aliquet
nisi arcu, sed vestibulum urna faucibus pretium. Maecenas laoreet diam
non urna tincidunt iaculis a ut ex. Aenean sem enim, laoreet id accumsan
sed, faucibus vitae diam. Aenean facilisis tincidunt risus. Mauris sit
amet hendrerit est, sit amet maximus augue.

\hypertarget{tables}{%
\section{Tables}\label{tables}}

You can use markdown tables, too. Captions are preceded by the
\emph{Table} keyword. Unfortunately, I can't figure out how to put the
caption below the table (it's above it by default).

\begin{longtable}[]{@{}cl@{}}
\caption{Table \label{param_table}\ref{param_table}:
\textbf{Compatibility flags}. A list of flags for compatibility testing
functions.}\tabularnewline
\toprule
Flag & Indication\tabularnewline
\midrule
\endfirsthead
\toprule
Flag & Indication\tabularnewline
\midrule
\endhead
1 & CONTENT\_ALL\_A\_IN\_B\tabularnewline
2 & CONTENT\_ALL\_B\_IN\_A\tabularnewline
4 & LENGTHS\_ALL\_A\_IN\_B\tabularnewline
8 & LENGTHS\_ALL\_B\_IN\_A\tabularnewline
16 & NAMES\_ALL\_A\_IN\_B\tabularnewline
32 & NAMES\_ALL\_B\_IN\_A\tabularnewline
64 & CONTENT\_A\_ORDER\tabularnewline
128 & CONTENT\_B\_ORDER\tabularnewline
256 & CONTENT\_ANY\_SHARED\tabularnewline
512 & LENGTHS\_ANY\_SHARED\tabularnewline
1024 & NAMES\_ANY\_SHARED\tabularnewline
\bottomrule
\end{longtable}

In hac habitasse platea dictumst. Mauris ut aliquet nunc, id mattis
velit. Nunc commodo enim sed orci ultrices sodales. Fusce nec sem est.
Nam euismod erat at neque facilisis iaculis. Cras rutrum elementum erat
eu egestas. In sit amet est vitae ligula semper vestibulum sit amet quis
justo. Sed porta dolor ac scelerisque congue.

Vivamus convallis arcu et lacus egestas tempus. Pellentesque habitant
morbi tristique senectus et netus et malesuada fames ac turpis egestas.
Integer eleifend efficitur risus sit amet suscipit. Curabitur
consectetur sapien eget nulla maximus, quis lobortis nisl porta. Sed
hendrerit semper placerat. Nulla sagittis orci arcu, a tincidunt lorem
lacinia sit amet. Nullam nec fringilla odio. In mollis vitae nibh ut
sollicitudin. Aliquam finibus tellus quis sollicitudin cursus. Cras
lobortis, tortor ac sodales tempus, lacus ipsum aliquam mauris, sed
placerat neque ante in arcu. Aliquam erat volutpat. Donec eu sodales
odio, eu cursus libero.

Sed in porttitor leo. Class aptent taciti sociosqu ad litora torquent
per conubia nostra, per inceptos himenaeos. Sed tristique malesuada
ligula, vulputate commodo mi lacinia ut. Ut consectetur, mauris vitae
hendrerit mattis, nisi urna condimentum nibh, vitae rhoncus nulla tellus
sit amet dui. Vivamus tempus magna eget quam posuere interdum. Nulla
turpis augue, consequat quis euismod a, tincidunt sed ligula. Integer ac
iaculis lacus, nec posuere augue. Fusce vitae dictum felis. In hac
habitasse platea dictumst. Etiam suscipit magna turpis, eget volutpat
lectus placerat quis. Mauris sed cursus erat. Sed a pellentesque felis.
Ut in blandit dolor, vitae lobortis justo. Aenean turpis felis, pulvinar
fringilla vulputate et, venenatis in lorem. Donec vulputate, nunc non
imperdiet ullamcorper, justo nunc placerat elit, ut pretium justo metus
a mi.

\hypertarget{embedded-latex}{%
\section{Embedded LaTeX}\label{embedded-latex}}

You can insert \texttt{latex} in-line in the markdown document:
\(rList[I_E] \leq q.start\)

Or you can create separate environments like this:

\hypertarget{algorithm-examples}{%
\subsection{Algorithm examples}\label{algorithm-examples}}

These examples use the \texttt{algorithmic} environment (from the
\texttt{algorithmcx} package:)

\begin{algorithmic}
\Require $n \geq 0$
\Ensure $y = x^n$
\State $y \Leftarrow 1$
\State $X \Leftarrow x$
\State $N \Leftarrow n$
\While{$N \neq 0$}
\If{$N$ is even}
  \State $X \Leftarrow X \times X$
  \State $N \Leftarrow \frac{N}{2} $  \Comment{This is a comment}
\ElsIf{$N$ is odd}
  \State $y \Leftarrow y \times X$
  \State $N \Leftarrow N - 1$
\EndIf
\EndWhile
\end{algorithmic}

\begin{algorithmic}[1]
\Repeat
\Comment{forever}
\State this\Until{you die.}
\end{algorithmic}

This example uses the \texttt{algorithm} environment:

\begin{algorithm}
\caption{Euclid’s algorithm}\label{euclid}
\begin{algorithmic}[1]
\Procedure{Euclid}{$a,b$}\Comment{The g.c.d. of a and b}
\State $r\gets a\bmod b$
\While{$r\not=0$}\Comment{We have the answer if r is 0}
\State $a\gets b$
\State $b\gets r$
\State $r\gets a\bmod b$
\EndWhile\label{euclidendwhile}
\State \textbf{return} $b$\Comment{The gcd is b}
\EndProcedure
\end{algorithmic}
\end{algorithm}

Maecenas vitae sodales est, venenatis ullamcorper magna. Integer id orci
ut arcu venenatis mattis. Pellentesque eget risus non lectus interdum
efficitur. In pharetra odio in tellus eleifend commodo. Morbi facilisis
mauris ac eros gravida pretium. Nam sit amet nisi massa. Morbi at turpis
in leo dictum suscipit. Ut interdum, orci sed laoreet venenatis, odio
dui consectetur tortor, sit amet vulputate ipsum neque ac ante. Vivamus
vitae mi interdum, dignissim leo lobortis, ultricies leo. Aenean
facilisis sagittis urna in blandit. Sed sit amet consectetur purus.
Mauris bibendum efficitur magna, vitae egestas lacus pretium dignissim.
Nullam eu magna est. Suspendisse vel lobortis metus.

Suspendisse potenti. Donec gravida ut mauris vel scelerisque. Nullam
gravida maximus porttitor. Duis dictum nisl sed neque tristique sodales.
Maecenas lacinia dolor eget ligula volutpat maximus. Etiam placerat
lobortis enim ut iaculis. Orci varius natoque penatibus et magnis dis
parturient montes, nascetur ridiculus mus. Fusce porta venenatis metus,
vehicula sagittis ligula faucibus vel. Nunc nibh ipsum, vulputate sed
bibendum non, euismod at sem. Praesent mi nulla, ornare vitae est a,
euismod facilisis mauris.

Mauris a orci vehicula, aliquam orci in, cursus eros. Lorem ipsum dolor
sit amet, consectetur adipiscing elit. Cras semper vel enim eu dapibus.
Maecenas placerat arcu nec metus tincidunt pharetra. Aenean rhoncus
lacinia elit et cursus. Ut et metus vel augue sagittis volutpat quis nec
nisi. Ut massa nisi, maximus vitae faucibus ut, eleifend ut odio.

Nam aliquam ex non accumsan efficitur. Nullam vehicula lorem vitae
porttitor pellentesque. Fusce a tristique mi, sed congue velit. Nullam
at ornare quam. Proin hendrerit accumsan ipsum, sed viverra velit
vehicula sit amet. Donec non lectus diam. Sed condimentum non velit vel
suscipit. Sed odio ex, vestibulum ullamcorper odio sit amet, lobortis
accumsan risus. Nulla facilisi. Mauris eleifend viverra metus, ac varius
lacus scelerisque non.

\hypertarget{references}{%
\section*{References}\label{references}}
\addcontentsline{toc}{section}{References}

\hypertarget{refs}{}
\begin{cslreferences}
\leavevmode\hypertarget{ref-Sheffield2016}{}%
1. Sheffield, N. C. \& Bock, C. LOLA: Enrichment analysis for genomic
region sets and regulatory elements in R and Bioconductor.
\emph{Bioinformatics} \textbf{32,} 587--589 (2016).

\leavevmode\hypertarget{ref-Sheffield2018}{}%
2. Sheffield, N. C., Nagraj, V. \& Reuter, V. simpleCache: R caching for
reproducible, distributed, large-scale projects. \emph{The Journal of
Open Source Software} \textbf{3,} 463 (2018).
\end{cslreferences}

% make the references smaller, and sloppier


\end{document}
